% Part III-A, discussion (included by main_en.tex).

\section{Discussion}\label{sec:discussion}

\paragraph{No sign that the visible tests were exploited.} The concern behind this paper is that a policy rewarded by
the tests it can see learns to satisfy those tests rather than to solve the problem. In our setting we saw no sign of
it within 100 steps. Training raised the share of problems solved on all three sets, including the two the policy
never trained on. The share of visible-test passes that are wrong fell by about half on KodCode validation and
HumanEval+ and barely moved on MBPP+. Programs that copy values from the visible tests stayed as rare as they were in
the base policy. The wrong programs that pass the visible tests are, before and after training, almost never copies of
the tests: they handle the cases the examples show and get others wrong. A visible-test reward pays for them, but the
base policy already wrote them, and on two of the three sets training made them a smaller share of the passes.

\paragraph{What the judge could do, and did.} Offline, both judges rank a problem's visible-test passes by whether
they also pass the held-out tests better than chance (\S\ref{sec:res-study}), and both catch hard-coded lookups. In
training, the judge factor did not change the outcome beyond the variation between seeds: no difference on KodCode
validation and HumanEval+, and on MBPP+ a gain in one seed and none in the other. What was left to suppress were not
copies of the tests but solutions that are wrong on cases the examples do not show. Whether a judge factor helps where
a policy does learn to special-case, with weaker tests, harder problems or longer training, these runs cannot say.

\paragraph{Better tests remove a few false passes, not unsolved problems.} The oracle reward, which pays only for programs
that pass every test, ended with fewer wrong programs passing the visible tests than \texttt{vis} on all three sets (by
an interval that excludes zero on MBPP+ only), and with the same share of problems solved.
Within 100 steps, knowing which visible-test passes are wrong did not teach the policy to solve more problems; it
taught it to pass the visible tests less often when it could not solve them. A perfect judge of correctness, used as
the factor, would give exactly this reward.

\paragraph{Limitations.} One policy size (4B), one task format (a single function, up to three attempts with the
first failing test as feedback), 100 steps, one learning rate taken from Part II, and two seeds per reward, which
disagree by up to \NUM{cgallsd_viskev_mbpp} points (standard deviation) on MBPP+. The intervals resample problems and not
seeds. The held-out tests are those of KodCode and EvalPlus, and a program that fails them is wrong by their
specification, which can be stricter than the problem statement. The special-casing check finds only programs that
contain a value from the visible tests; special-casing without such a value (a branch on the input's length, say) is
not detected. Other experiments shared the GPUs during training.
